\chapter{Symmetry and Conservation Laws}

\lectureoverview{Stationary action and the Euler--Lagrange equations are one
half of the heart of classical mechanics; the Hamiltonian formulation is the
other.  Here we develop the first half far enough to expose a deep pattern:
every continuous symmetry of a Lagrangian produces a conserved quantity.  We
shall discover the pattern in translations and rotations, prove it for a
general infinitesimal transformation, and then examine carefully where the
argument does and does not apply.}

\section{The Minimum Mathematics That Explains the Physics}

The marks already on the board may look a little like the Schrödinger
equation, but that is not tonight's subject.  The subject is symmetry and
conservation.  There is no honest way to explain their connection without
equations.  The equations are elementary; what matters is arranging them so
that the structure becomes visible.  Scientific American may name the result,
but here we want the minimum mathematics required to derive it as physicists
would.

Introduce a small parameter \(\epsilon\).  We retain terms proportional to
\(\epsilon\) and discard \(O(\epsilon^2)\).  This is shorthand for a limiting
argument, not a claim that a nonzero real number literally has zero square.
For an ordinary function \(f(x,y)\),
\begin{equation}
  \delta f
  =\frac{\partial f}{\partial x}\,\delta x
   +\frac{\partial f}{\partial y}\,\delta y
   +O(\epsilon^2),
  \label{eq:first-variation-f}
\end{equation}
where both \(\delta x\) and \(\delta y\) are proportional to \(\epsilon\).
The symbol \(\delta f\) means the change in \(f\), not the product of a number
called \(\delta\) with \(f\).  A stationary point obeys
\begin{equation}
  \delta f=0
  \qquad\hbox{to first order.}
  \label{eq:stationary-f}
\end{equation}

For a mechanical path \(q_i(t)\), the analogous quantity is the action
\begin{equation}
  S[q]=\int_{t_1}^{t_2}L(q,\dot q,t)\,dt.
  \label{eq:action-l4}
\end{equation}
Varying the path while fixing its endpoints gives
\begin{equation}
  \delta S
  =\int_{t_1}^{t_2}dt\,
   \sum_i\left(
     \frac{\partial L}{\partial q_i}\,\delta q_i
     +\frac{\partial L}{\partial\dot q_i}\,\delta\dot q_i
   \right).
  \label{eq:action-variation-l4}
\end{equation}
The two terms arise because changing a path changes both its position and its
velocity.  Integration by parts, together with
\(\delta q_i(t_1)=\delta q_i(t_2)=0\), turns \(\delta S=0\) into
\begin{equation}
  \frac{d}{dt}\left(\frac{\partial L}{\partial\dot q_i}\right)
  -\frac{\partial L}{\partial q_i}=0.
  \label{eq:el-l4}
\end{equation}

\begin{figure}[htbp]
  \centering
  \includegraphics[width=0.76\textwidth]{lecture_04_figure_05.png}
  \caption{The Euler--Lagrange equation recalled as the starting point for
  symmetry and conservation.}
  \lectureframe{00:09:00}
\end{figure}

The correct phrase is \emph{stationary action}, not necessarily least action.
That distinction will return near the end.

\section{Generalized Coordinates and Conjugate Momentum}

Cartesian positions are too narrow a language for general mechanics.  The
coordinates may be particle positions, angles, collective variables, or field
amplitudes, so we call them \(q_i\).  A Lagrangian is a compressed
specification of the dynamics.  Even the Standard Model, written out without
very condensed notation, occupies roughly a page; writing a separate equation
of motion for every field component would occupy vastly more.  The one
function \(L\) packages all of them.

The momentum conjugate to \(q_i\) is
\begin{equation}
  p_i=\frac{\partial L}{\partial\dot q_i}.
  \label{eq:canonical-momentum-l4}
\end{equation}
This is why generalized coordinates are traditionally called \(q\)'s: the
corresponding momenta are \(p\)'s, giving us our \(p\)'s and \(q\)'s.  In this
notation Eq.~\eqref{eq:el-l4} becomes
\begin{equation}
  \dot p_i=\frac{\partial L}{\partial q_i}.
  \label{eq:momentum-el-l4}
\end{equation}

For the familiar one-dimensional particle,
\begin{equation}
  L=\frac{m\dot x^2}{2}-V(x),
  \qquad
  \frac{\partial L}{\partial\dot x}=m\dot x,
  \label{eq:particle-momentum-l4}
\end{equation}
so the conjugate momentum is ordinary momentum.

\begin{figure}[htbp]
  \centering
  \includegraphics[width=0.72\textwidth]{lecture_04_figure_02.png}
  \caption{The kinetic-minus-potential Lagrangian used to identify conjugate
  momentum with \(m\dot x\).}
  \lectureframe{00:14:48}
\end{figure}

If \(q_i\) does not appear in \(L\), then \(\dot p_i=0\).  This is the familiar
cyclic-coordinate route to a conservation law.  The real connection is more
general, and it is useful to discover it in examples before proving it.

\section{Conservation Laws Found Before Their Explanation}

\subsection{Two coordinates and a separation potential}

Set both masses equal to one and take
\begin{equation}
  L=\frac{\dot q_1^{\,2}+\dot q_2^{\,2}}{2}-V(q_1-q_2).
  \label{eq:separation-lagrangian}
\end{equation}
One may picture two particles on a line.  Their interaction depends only on
their separation.

\begin{figure}[htbp]
  \centering
  \includegraphics[width=0.78\textwidth]{lecture_04_figure_03.png}
  \caption{The two-coordinate Lagrangian and the beginning of its first
  momentum equation.}
  \lectureframe{00:18:56}
\end{figure}

Writing \(V'\) for differentiation with respect to the argument of \(V\),
\begin{align}
  \dot p_1&=-V'(q_1-q_2),
  \notag\\
  \dot p_2&=+V'(q_1-q_2).
  \label{eq:separation-momentum-equations}
\end{align}
The second sign reverses because increasing \(q_2\) decreases \(q_1-q_2\).

\begin{classroomqa}
  \audiencequestion{Should the first momentum equation contain \(V'\), and
  where does its minus sign come from?
  \lecturetimestamp{00:20:52--00:21:35}}
  \lecturerresponse{The Euler--Lagrange equation says
  \(\dot p_1=\partial L/\partial q_1\), not that the Lagrangian is zero.
  Since \(L=T-V\), differentiating with respect to \(q_1\) gives
  \(-V'\).  Differentiating with respect to \(q_2\) supplies one more minus
  sign from \(q_1-q_2\), so the second equation has \(+V'\).}
\end{classroomqa}

\begin{figure}[htbp]
  \centering
  \includegraphics[width=0.78\textwidth]{lecture_04_figure_04.png}
  \caption{The sign-sensitive derivative in the \(p_1\) equation; the
  \(p_2\) equation receives the opposite sign.}
  \lectureframe{00:20:41}
\end{figure}

Adding Eqs.~\eqref{eq:separation-momentum-equations} gives
\begin{equation}
  \frac{d}{dt}(p_1+p_2)=0.
  \label{eq:total-p-conserved}
\end{equation}
We have found a conserved quantity, but have not yet explained why it exists.

\subsection{A less obvious linear combination}

Now let
\begin{equation}
  L=\frac{\dot q_1^{\,2}+\dot q_2^{\,2}}{2}
    -V(Aq_1+Bq_2),
  \label{eq:ab-lagrangian}
\end{equation}
where \(A\) and \(B\) are constants.  Then
\begin{align}
  \dot p_1&=-A\,V'(Aq_1+Bq_2),
  \notag\\
  \dot p_2&=-B\,V'(Aq_1+Bq_2).
  \label{eq:ab-momentum-equations}
\end{align}
Multiplying the first equation by \(B\), the second by \(A\), and subtracting
eliminates the force:
\begin{equation}
  \frac{d}{dt}(Bp_1-Ap_2)=0.
  \label{eq:ab-conserved}
\end{equation}

\begin{figure}[htbp]
  \centering
  \includegraphics[width=0.78\textwidth]{lecture_04_figure_06.png}
  \caption{The weighted momentum equations whose force terms cancel in
  \(Bp_1-Ap_2\).}
  \lectureframe{00:24:20}
\end{figure}

\begin{classroomqa}
  \audiencequestion{Where might constants such as \(A\) and \(B\) come from
  physically?\lecturetimestamp{00:25:09--00:30:15}}
  \lecturerresponse{Consider two particles of masses \(m\) and \(M\):
  \[
    L=\frac{m\dot X^2}{2}+\frac{M\dot Y^2}{2}-V(X-Y).
  \]
  Rescale to \(Q_1=\sqrt m\,X\) and \(Q_2=\sqrt M\,Y\).  The kinetic term
  becomes \((\dot Q_1^2+\dot Q_2^2)/2\), while the potential depends on
  \(Q_1/\sqrt m-Q_2/\sqrt M\).  Thus unequal masses naturally become
  coefficients in a linear combination.  Replacing \(Q_2\) by \(-Q_2\)
  changes the displayed minus to a plus without changing its squared kinetic
  term.}
\end{classroomqa}

\begin{figure}[htbp]
  \centering
  \includegraphics[width=0.82\textwidth]{lecture_04_figure_07.png}
  \caption{Mass-weighted coordinates simplify the kinetic energy while
  moving the mass factors into the potential.}
  \lectureframe{00:28:40}
\end{figure}

\begin{classroomqa}
  \audiencequestion{Does this show that a conserved quantity can exist
  without a cyclic coordinate?\lecturetimestamp{00:31:13--00:31:58}}
  \lecturerresponse{There is no \emph{obvious} cyclic coordinate in the
  variables as written.  A change of variables may reveal one, but hunting
  for it is not the useful viewpoint here.  The symmetry argument will
  identify the conserved combination directly.}
\end{classroomqa}

\section{What We Mean by a Symmetry}

\subsection{Translations}

Begin with a free particle of unit mass,
\begin{equation}
  L=\frac{\dot q^2}{2}.
  \label{eq:free-unit-l4}
\end{equation}
Translate every point of its history by the same constant:
\begin{equation}
  q'(t)=q(t)+\epsilon.
  \label{eq:free-translation}
\end{equation}
Since \(\epsilon\) is time independent, \(\dot q'=\dot q\), and therefore
\(\delta L=0\).  This invariance is translation symmetry.

There are two equivalent pictures.  Passively, shift the coordinate origin
while leaving the object fixed.  Actively, keep the axes fixed and move the
object.  The active language will prove less confusing: move the physical
configuration and ask whether \(L\) notices.

For Eq.~\eqref{eq:separation-lagrangian}, translate both coordinates:
\begin{equation}
  \delta q_1=\epsilon,
  \qquad
  \delta q_2=\epsilon.
  \label{eq:common-translation}
\end{equation}
The velocities are unchanged and so is \(q_1-q_2\).

\begin{figure}[htbp]
  \centering
  \includegraphics[width=0.78\textwidth]{lecture_04_figure_08.png}
  \caption{A common translation leaves the separation of two coordinates
  unchanged.}
  \lectureframe{00:35:30}
\end{figure}

For Eq.~\eqref{eq:ab-lagrangian}, use the less obvious transformation
\begin{equation}
  \delta q_1=B\epsilon,
  \qquad
  \delta q_2=-A\epsilon.
  \label{eq:ab-transformation}
\end{equation}
Again the velocities do not change, while
\begin{equation}
  \delta(Aq_1+Bq_2)=AB\epsilon-BA\epsilon=0.
  \label{eq:ab-invariant-argument}
\end{equation}
Thus this transformation is a symmetry.

\begin{figure}[htbp]
  \centering
  \includegraphics[width=0.80\textwidth]{lecture_04_figure_09.png}
  \caption{The asymmetric translation that preserves \(Aq_1+Bq_2\).}
  \lectureframe{00:40:20}
\end{figure}

\begin{classroomqa}
  \audiencequestion{Do the coefficients encode the masses, so that the
  conserved quantity becomes ordinary center-of-mass momentum?
  \lecturetimestamp{00:42:25--00:42:45}}
  \lecturerresponse{Yes.  In the mass-weighted variables the symmetry and
  conserved combination look asymmetric.  Transforming back to \(X\) and
  \(Y\) recovers the ordinary simultaneous translation and total momentum.}
\end{classroomqa}

\subsection{Rotations}

Consider one particle in the plane,
\begin{equation}
  L=\frac{m}{2}(\dot x^2+\dot y^2)-V(x^2+y^2).
  \label{eq:central-lagrangian-l4}
\end{equation}
This is a central potential.  We compare two fixed coordinate systems rotated
relative to one another; the axes are not rotating in time.  With one
consistent orientation,
\begin{align}
  x'&=x\cos\theta+y\sin\theta,
  \notag\\
  y'&=-x\sin\theta+y\cos\theta.
  \label{eq:finite-rotation-l4}
\end{align}

\begin{figure}[htbp]
  \centering
  \includegraphics[width=0.82\textwidth]{lecture_04_figure_10.png}
  \caption{A central-potential Lagrangian and two fixed Cartesian frames
  related by a rotation.}
  \lectureframe{00:45:40}
\end{figure}

Set \(\theta=\epsilon\) and keep first order:
\begin{equation}
  \cos\epsilon=1+O(\epsilon^2),
  \qquad
  \sin\epsilon=\epsilon+O(\epsilon^3).
  \label{eq:small-angle-l4}
\end{equation}
The infinitesimal rotation is
\begin{align}
  \delta x&=y\epsilon,
  &\delta y&=-x\epsilon,
  \notag\\
  \delta\dot x&=\dot y\epsilon,
  &\delta\dot y&=-\dot x\epsilon.
  \label{eq:infinitesimal-rotation-l4}
\end{align}

\begin{figure}[htbp]
  \centering
  \includegraphics[width=0.82\textwidth]{lecture_04_figure_11.png}
  \caption{The infinitesimal rotation of coordinates and velocities.}
  \lectureframe{00:49:50}
\end{figure}

Unlike a translation, the amount by which \(x\) changes depends on \(y\), and
the amount by which \(y\) changes depends on \(x\).  Nevertheless,
\begin{align}
  \delta(x^2+y^2)
  &=2x\,\delta x+2y\,\delta y=0,
  \notag\\
  \delta(\dot x^2+\dot y^2)
  &=2\dot x\,\delta\dot x+2\dot y\,\delta\dot y=0.
  \label{eq:rotation-invariants-l4}
\end{align}
Both the potential and kinetic terms are invariant.

\begin{figure}[htbp]
  \centering
  \includegraphics[width=0.76\textwidth]{lecture_04_figure_12.png}
  \caption{The first-order cancellation proving that the squared radius is
  rotationally invariant.}
  \lectureframe{00:53:20}
\end{figure}

\section{The General Conservation Theorem}

The translation and rotation examples have one form in common.  Let a
continuous infinitesimal transformation be
\begin{equation}
  \delta q_i=\epsilon f_i(q),
  \qquad
  \delta\dot q_i=\frac{d}{dt}\delta q_i,
  \label{eq:general-infinitesimal-symmetry}
\end{equation}
where \(\epsilon\) is time independent and \(f_i\) may depend on all the
coordinates.

\begin{figure}[htbp]
  \centering
  \includegraphics[width=0.78\textwidth]{lecture_04_figure_13.png}
  \caption{A general infinitesimal displacement field
  \(\delta q_i=\epsilon f_i(q)\).}
  \lectureframe{00:56:00}
\end{figure}

The first-order variation of the Lagrangian is
\begin{equation}
  \delta L
  =\sum_i\left(
    \frac{\partial L}{\partial q_i}\delta q_i
    +\frac{\partial L}{\partial\dot q_i}\delta\dot q_i
  \right).
  \label{eq:delta-l-general}
\end{equation}

\begin{figure}[htbp]
  \centering
  \includegraphics[width=0.84\textwidth]{lecture_04_figure_14.png}
  \caption{The two contributions to \(\delta L\): coordinate variation and
  velocity variation.}
  \lectureframe{01:02:30}
\end{figure}

Now restrict attention to a physical path, so the Euler--Lagrange equations
hold.  Using
\begin{equation}
  \frac{\partial L}{\partial\dot q_i}=p_i,
  \qquad
  \frac{\partial L}{\partial q_i}=\dot p_i,
  \label{eq:el-substitution-l4}
\end{equation}
we find
\begin{align}
  \delta L
  &=\sum_i\left(\dot p_i\,\delta q_i
                +p_i\,\delta\dot q_i\right)
  \notag\\
  &=\frac{d}{dt}\sum_i p_i\delta q_i
  =\epsilon\frac{d}{dt}\sum_i p_i f_i(q).
  \label{eq:noether-derivation-l4}
\end{align}
The central step is only the product rule.

If the transformation is a symmetry, \(\delta L=0\).  Since \(\epsilon\) is
an arbitrary constant factor,
\begin{equation}
  \boxed{
  \frac{d\mathcal Q}{dt}=0,
  \qquad
  \mathcal Q=\sum_i p_i f_i(q)}
  .
  \label{eq:noether-charge-l4}
\end{equation}
This is the classical-mechanics core of Noether's theorem for exact,
time-independent point symmetries.

\begin{figure}[htbp]
  \centering
  \includegraphics[width=0.84\textwidth]{lecture_04_figure_15.png}
  \caption{The product-rule step that turns invariance of \(L\) into a
  conserved quantity.}
  \lectureframe{01:05:00}
\end{figure}

\begin{classroomqa}
  \audiencequestion{In \(f_i(q)\), does \(q\) carry another index?
  \lecturetimestamp{01:06:57--01:07:31}}
  \lecturerresponse{The unindexed \(q\) denotes the whole set
  \((q_1,q_2,\ldots)\).  Each component \(f_i\) may depend on all those
  coordinates.  For rotation, for example, \(f_x=y\) and \(f_y=-x\).}
\end{classroomqa}

The examples now cease to be accidents.  A common translation has
\(f_1=f_2=1\), so
\begin{equation}
  \mathcal Q=p_1+p_2.
  \label{eq:noether-common-p}
\end{equation}
The asymmetric translation has \(f_1=B\), \(f_2=-A\), so
\begin{equation}
  \mathcal Q=Bp_1-Ap_2.
  \label{eq:noether-ab-p}
\end{equation}
For rotation, \(f_x=y\) and \(f_y=-x\):
\begin{equation}
  \mathcal Q=p_xy-p_yx.
  \label{eq:rotation-charge-sign}
\end{equation}
Up to the orientation sign, this is angular momentum,
\begin{equation}
  J=xp_y-yp_x.
  \label{eq:angular-momentum-l4}
\end{equation}

\begin{figure}[htbp]
  \centering
  \includegraphics[width=0.82\textwidth]{lecture_04_figure_16.png}
  \caption{The rotational Noether charge, equal up to orientation to planar
  angular momentum.}
  \lectureframe{01:11:00}
\end{figure}

For many interacting particles, rotate every particle together.  Invariance
then conserves the sum of their angular momenta, not each particle's angular
momentum separately.

\section{Questions About the Theorem's Hypotheses}

\begin{classroomqa}
  \audiencequestion{Would a rotating set of axes introduce explicit time
  dependence into \(\delta L\)?\lecturetimestamp{01:12:57--01:13:57}}
  \lecturerresponse{Yes.  The proof just given assumes a time-independent
  symmetry parameter: a fixed translation or a fixed rotation.  A
  transformation whose parameter depends on time produces additional terms
  and requires a more general treatment.}
\end{classroomqa}

\begin{classroomqa}
  \audiencequestion{Is the capital \(Q\) in the board formula the
  transformation itself?\lecturetimestamp{01:14:02--01:15:58}}
  \lecturerresponse{No.  That notation unfortunately collides with the
  coordinates \(q_i\).  The transformation is
  \(\delta q_i=\epsilon f_i(q)\); the resulting conserved quantity is
  \(\mathcal Q=\sum_i p_i f_i\).  The coefficient \(f_i\) tells how the
  \(i\)-th coordinate moves under the symmetry.}
\end{classroomqa}

Energy conservation is also generated by a symmetry, but a different one:
translation in time, or independence of where the zero of time is placed.
That case and the Hamiltonian formulation are deferred.

\begin{classroomqa}
  \audiencequestion{Could the actual motion be absorbed into a changing
  coordinate system, so mechanics becomes nothing but coordinate
  transformations?\lecturetimestamp{01:17:43--01:20:01}}
  \lecturerresponse{One can choose coordinates that follow a particular
  motion, but then the transformation depends on the initial state and on the
  trajectory.  The symmetry transformation used here is fixed independently
  of which solution the system follows.  Motion itself is more naturally
  represented in the combined space of coordinates and momenta.}
\end{classroomqa}

\section{The Harmonic Oscillator in Phase Space}

With mass and spring constant set to one,
\begin{equation}
  L=\frac{\dot q^2}{2}-\frac{q^2}{2}.
  \label{eq:oscillator-l4}
\end{equation}
The conjugate momentum and equation of motion give
\begin{equation}
  p=\dot q,
  \qquad
  \dot p=-q,
  \label{eq:oscillator-first-order-l4}
\end{equation}
or
\begin{equation}
  \dot q=p,
  \qquad
  \dot p=-q.
  \label{eq:oscillator-phase-flow}
\end{equation}
These equations describe uniform circular motion in the \(q\)-\(p\) plane.
The back-and-forth real-space oscillator is the projection of that circle onto
the \(q\)-axis: speed is largest at \(q=0\) and vanishes at the turning
points.

\begin{figure}[htbp]
  \centering
  \includegraphics[width=0.82\textwidth]{lecture_04_figure_17.png}
  \caption{The harmonic-oscillator equations and their circular orbit in
  \(q\)-\(p\) space.}
  \lectureframe{01:25:30}
\end{figure}

\begin{classroomqa}
  \audiencequestion{Does the phase-space circle run clockwise?
  \lecturetimestamp{01:25:47--01:27:15}}
  \lecturerresponse{With \(p\) horizontal and \(q\) vertical, yes.  At
  maximal positive \(q\), \(\dot p=-q<0\), so the motion points downward.
  Reversing the placement of the axes reverses the visual sense, which
  explains why the convention is easy to remember incorrectly.}
\end{classroomqa}

The pair in Eq.~\eqref{eq:oscillator-phase-flow} already resembles Hamilton's
equations, but it has been obtained entirely within Lagrangian mechanics.

\section{Continuous and Discrete Symmetries}

\begin{classroomqa}
  \audiencequestion{Must a symmetry be expressed by an infinitesimal
  Cartesian transformation?\lecturetimestamp{01:27:48--01:30:31}}
  \lecturerresponse{Cartesian coordinates are only convenient; the essential
  requirement for this proof is continuity.  Some symmetries have no
  infinitesimal version.  For an even potential, reflection sends
  \(q\mapsto -q\) and \(\dot q\mapsto-\dot q\), leaving \(L\) unchanged, but
  a reflection cannot be assembled from arbitrarily small reflections.
  This classical derivation therefore produces no associated conserved
  quantity.}
\end{classroomqa}

\begin{figure}[htbp]
  \centering
  \includegraphics[width=0.82\textwidth]{lecture_04_figure_18.png}
  \caption{An even potential is invariant under the discrete reflection
  \(q\mapsto -q\).}
  \lectureframe{01:29:30}
\end{figure}

\begin{classroomqa}
  \audiencequestion{If the conserved quantity \(\mathcal Q\) is nonzero, does
  that show the transformation was not a symmetry?
  \lecturetimestamp{01:30:40--01:31:30}}
  \lecturerresponse{No.  A conserved quantity may have any constant value.
  The relevant test is \(\dot{\mathcal Q}\).  If the quantity derived from a
  proposed transformation is not conserved, then that proposed
  transformation was not a symmetry under the hypotheses of the proof.}
\end{classroomqa}

\begin{classroomqa}
  \audiencequestion{Reflection still leaves the Lagrangian unchanged, so why
  is there no conservation law from it?\lecturetimestamp{01:31:34--01:32:27}}
  \lecturerresponse{Invariance alone is visible, but the derivation passed
  through a first-order displacement and the product rule.  Reflection moves
  one side of configuration space to the other in a single discrete step.
  With no infinitesimal generator \(f_i\), the formula
  \(\mathcal Q=\sum_i p_i f_i\) has nothing to use.}
\end{classroomqa}

\begin{classroomqa}
  \audiencequestion{What is the name for the property of being built from
  arbitrarily small transformations?\lecturetimestamp{01:32:36--01:33:23}}
  \lecturerresponse{It is a \emph{continuous symmetry}.  Translation and
  rotation are continuous; reflection is a \emph{discrete symmetry}.}
\end{classroomqa}

\begin{classroomqa}
  \audiencequestion{If many infinitesimal steps make a finite
  transformation, can the discarded \(O(\epsilon^2)\) errors accumulate?
  \lecturetimestamp{01:33:23--01:36:10}}
  \lecturerresponse{The infinitesimal invariance must hold at every point
  reached along the transformation, not only at the starting point.  Divide a
  finite parameter interval into \(N\) steps of size \(\epsilon\): the
  accumulated first-order displacement is \(N\epsilon\), while a local
  \(O(\epsilon^2)\) remainder accumulates as
  \(N\epsilon^2=(N\epsilon)\epsilon\to0\).  This is the local-to-finite
  structure formalized by continuous transformation groups.}
\end{classroomqa}

\begin{classroomqa}
  \audiencequestion{Are discrete symmetries physically realized in classical
  systems?\lecturetimestamp{01:36:32--01:37:25}}
  \lecturerresponse{Certainly.  The harmonic oscillator itself has the even
  potential \(V(q)\propto q^2\), so it is invariant under
  \(q\mapsto -q\).  There are many such examples; the point is only that this
  discrete invariance supplies no new classical Noether charge.}
\end{classroomqa}

\begin{classroomqa}
  \audiencequestion{But if both continuous and discrete transformations leave
  \(L\) unchanged, is that not the same condition?
  \lecturetimestamp{01:37:37--01:39:18}}
  \lecturerresponse{It is the same finite statement of invariance, but not the
  same mathematical input to this proof.  The proof needs a first-order
  variation.  For a reflection, 'first order' has no meaning, so the chain of
  equations leading to a charge cannot even be started.}
\end{classroomqa}

\begin{classroomqa}
  \audiencequestion{How does this relate historically to Noether's theorem?
  \lecturetimestamp{01:39:26--01:40:42}}
  \lecturerresponse{This is a restricted part of Emmy Noether's much more
  general theorem.  Particular links between symmetries and conserved
  quantities were known earlier, but Noether placed them in a unified and
  remarkably general framework in the early twentieth century.  Her
  mathematical importance was not matched by fair academic treatment.}
\end{classroomqa}

\begin{classroomqa}
  \audiencequestion{Is classical mechanics still an active research subject?
  \lecturetimestamp{01:40:47--01:42:37}}
  \lecturerresponse{Yes.  The elementary theorem used here is old, but
  nonlinear dynamics, integrable systems, chaos, fluids, and geometric
  mechanics remain rich subjects.  The symmetry principle also continues far
  beyond classical particle mechanics.}
\end{classroomqa}

\section{Why the Action Principle Matters}

\begin{classroomqa}
  \audiencequestion{Is there an intuitive reason that symmetry should imply
  conservation?\lecturetimestamp{01:42:49--01:47:57}}
  \lecturerresponse{Not from symmetry alone.  Consider viscous drag,
  \[
    m\ddot x=-c\dot x.
  \]
  The equation is invariant under \(x\mapsto x+\epsilon\), yet with
  \(p=m\dot x\),
  \[
    \dot p=-\frac{c}{m}p,
    \qquad
    p(t)=p(0)e^{-ct/m},
  \]
  so momentum decays.  This is an effective equation for a subsystem, not a
  fundamental stationary-action law for \(x\) alone.  Include the fluid
  molecules and the rest of the closed system, and total momentum is
  conserved.  Symmetry and conservation are tied together through the
  Lagrangian action principle, not through the appearance of a differential
  equation by itself.}
\end{classroomqa}

\begin{figure}[htbp]
  \centering
  \includegraphics[width=0.80\textwidth]{lecture_04_figure_19.png}
  \caption{Translation-invariant viscous drag, whose reduced particle
  momentum nevertheless decays.}
  \lectureframe{01:44:30}
\end{figure}

\begin{classroomqa}
  \audiencequestion{Can we quantify when the physical path is not a minimum
  of the action?\lecturetimestamp{01:48:08--01:51:38}}
  \lecturerresponse{A true minimum is exceptional; the action is generally a
  saddle in path space.  For the unit oscillator,
  \[
    S[q]=\frac12\int(\dot q^2-q^2)\,dt.
  \]
  A narrow, rapidly changing deformation raises the positive
  \(\dot q^2\) contribution.  A broad deformation that remains slightly
  displaced for a long time can make the negative \(-q^2\) contribution
  dominate and lower the action.  The same stationary path therefore has
  increasing directions and decreasing directions.  Quantum mechanics,
  through the path-integral viewpoint, gives the deeper origin of
  stationarity, but that argument is deferred.}
\end{classroomqa}

\begin{figure}[htbp]
  \centering
  \includegraphics[width=0.80\textwidth]{lecture_04_figure_20.png}
  \caption{A broad neighboring oscillator path illustrates why stationary
  action is typically a saddle rather than a minimum.}
  \lectureframe{01:50:20}
\end{figure}

\begin{classroomqa}
  \audiencequestion{Is there a one-to-one correspondence between continuous
  symmetries and conservation laws?\lecturetimestamp{01:51:44--01:52:43}}
  \lecturerresponse{Within the variational setting and after identifying
  equivalent or trivial cases, that is essentially the right picture.  Every
  continuous symmetry has a charge, and conserved quantities can generally be
  related back to symmetries.  Precise formulations carry technical
  qualifications, so the broad statement should not be used carelessly.}
\end{classroomqa}

\begin{classroomqa}
  \audiencequestion{Must every Lagrangian possess a nontrivial symmetry?
  \lecturetimestamp{01:52:43--01:53:23}}
  \lecturerresponse{No.  The examples were chosen to have symmetries.  A
  generic Lagrangian need not be invariant under any proposed spatial
  transformation.}
\end{classroomqa}

\begin{classroomqa}
  \audiencequestion{Should a symmetry finally be understood as an active
  motion or a passive change of coordinates?
  \lecturetimestamp{01:53:25--01:55:08}}
  \lecturerresponse{Either description can be made equivalent, but the active
  one is cleaner here.  Move the whole physical system and ask whether the
  Lagrangian changes.  In passive language one must instead say that the
  \emph{form} of the Lagrangian is unchanged when the coordinate frame moves.}
\end{classroomqa}

\begin{classroomqa}
  \audiencequestion{What comes next?\lecturetimestamp{01:55:10--01:55:35}}
  \lecturerresponse{Conservation of energy, its origin in time-translation
  symmetry, and then the Hamiltonian formulation.}
\end{classroomqa}
